\documentclass[sigconf,anonymous,review]{acmart}

\settopmatter{printacmref=true,printfolios=true}
\setcopyright{none}
\acmConference[ICAIF '26]{The 7th ACM International Conference on AI in
Finance}{November 14--17, 2026}{Milan, Italy}
\acmYear{2026}
\acmDOI{}
\acmISBN{}
\renewcommand\footnotetextcopyrightpermission[1]{}
\setlength{\emergencystretch}{1em}

\title{When Should an Economic Agent Trust Its Memory? Evidence-Grounded Rule
Evolution under Endogenous Feedback}

\begin{document}

\begin{abstract}
Long-horizon language-model agents often turn a small number of action--outcome
experiences into reusable natural-language rules. A biased reflection can then
become a persistent instruction, alter later actions, and change the environment
that generates subsequent evidence. We formulate this as a rule-reliability
problem using Evidence-Grounded Rule Memory (EGRM), a source-pinned lifecycle
implementation inherited from FinEvo. We do not claim that lifecycle code as a
new implementation contribution. The candidate contribution is the reliability
formulation plus an oracle-labeled switching benchmark, registered policy
comparisons, controlled false rules, and matched checkpoint continuations. The
primary outcomes are unsupported activation, harmful exposure, retirement
latency, and paired downstream utility loss. This pre-experiment manuscript
defines the protocol; effectiveness results are not yet available.
\end{abstract}

\maketitle

\section{Introduction}
Persistent memory lets an economic agent reuse experience across long horizons,
but it also creates a closed-loop failure mode: the agent writes a rule from its
own outcomes, follows that rule later, and thereby changes the evidence used to
judge it. Retrieval quality alone cannot answer when such a rule deserves
authority. We ask when a self-generated rule should be admitted, activated,
revised, or retired.

The evaluated rule lifecycle is inherited from a source-pinned FinEvo
implementation; this paper does not claim its code or module structure as a new
implementation contribution. Our candidate contribution is threefold. First,
we formalize rule reliability under endogenous feedback. Second, an
oracle-labeled regime-switch benchmark exposes known validity and
counterevidence. Third, predeclared verifier comparisons, a forced-active test,
and hash-bound continuations hold the environment, model, prompt, decoding, and
random state fixed across memory treatments. These
remain design contributions until the protocol is completed, preregistered, and
executed.

\section{Related Work and Scope}
Economic LLM agents model heterogeneous economic decisions and retain market
experience~\cite{li2024econagent}, while recent work studies
sentiment, memory weighting, and dynamic economic adaptation~\cite{liu2026econai}.
Experience-based agents store verbal reflections~\cite{shinn2023reflexion},
extract reusable insights~\cite{zhao2024expel}, and learn instructions from
past reflections~\cite{gupta2024metareflection}. Recent empirical work studies
experience following and its errors, including misaligned replay and feedback
from later evaluations~\cite{xiong2026memory}. These works rule out any claim
that EGRM is the first system for reflection from experience, memory of rules,
or feedback-based regulation.

General agent-memory research further addresses dynamic
organization~\cite{xu2025amem}, incremental memory
evaluation~\cite{hu2025memoryagentbench}, and safety under contaminated
memories~\cite{wei2025amemguard,xie2026memevobench}. The candidate distinction
here is narrower: reliability of self-generated action--outcome rules when rule
use changes the endogenous environment producing later evidence, evaluated
with known truth and matched interventions. This remains a scoped hypothesis,
not a universal novelty claim.

\section{Problem Setup}
At decision time $t$, agent $i$ observes state $s_{i,t}$, retrieves finalized
episodes, and executes action $a_{i,t}$. Only after the environment transition
is an episode $e_{i,t}=(s_{i,t},a_{i,t},s_{i,t+1},u_{i,t+1})$ eligible as
evidence. A candidate rule contains a context scope, state predicate, action
guidance, rationale, and cited episode identifiers. The verifier fixes the
outcome criterion before proposal generation.

\section{Evidence-Grounded Rule Memory}
The lifecycle evaluated here is inherited from the source-pinned FinEvo
implementation. The description below defines the evaluation subject and its
registered controls; it is not a claim that EGRM newly implements these
lifecycle mechanisms.

For each in-scope episode, EGRM records one of five mutually exclusive classes:
support, harmful compliance, alternative success, alternative failure, or
irrelevant. The default positive and negative scores are
\begin{equation}
 S= n_{\mathrm{sup}},\qquad
 C=n_{\mathrm{harm}}+0.5n_{\mathrm{alt\mbox{-}succ}},
\end{equation}
with margin $S-C$ and smoothed confidence
$(S+1)/(S+C+2)$. Alternative failure and irrelevant observations remain in the
audit denominator with zero evidentiary weight.

A proposal can be rejected or provisional, never directly active. Activation
requires distinct post-proposal support plus frozen margin and confidence
thresholds. Repeated harmful compliance or confidence decay retires a rule.
New versions must identify the superseded rule and cite genuinely new support.
An unverified-immediate control uses the same structured rule language but
bypasses evidence admission and retirement.

\section{Evaluation Protocol}
The controlled benchmark contains oracle-labeled rule families and regime
switches. The closed-loop study uses fresh paired seeds and predeclares verifier
on/off $\times$ error on/off contrasts, candidate-admission tests, and
forced-active false-rule treatments. From a shared checkpoint, matched A/A and
error branches with and without verification continue all agents for six
periods with identical random states.

A fixed false-rule injection tests how the registered policies respond to a
controlled known error. It does not estimate the frequency or character of
naturally hallucinated LLM proposals, and it cannot by itself establish that the
method mitigates such hallucinations. That broader claim would require a
separately registered natural-proposal audit.

The seed is the primary inference unit. Parse, provider, integrity, and budget
failures remain in the intention-to-treat denominator; failed seeds are not
replaced. A downstream effect must exceed the maximum matched A/A null, exceed
one discrete action bin, and agree in at least four of five seeds. Otherwise an
action-only change is described as prompt sensitivity.

\section{Results}
\textbf{Result status: not run.} The final version will report every registered
cell, raw paired deltas, direction counts, mean, median, and range. Historical
macroeconomic tables are excluded from EGRM results. If the evidence-aware arm
does not reduce unsupported activation, harmful exposure, and paired utility
loss under the registered go rule, the corresponding effectiveness claim will
be withdrawn.

\section{Limitations and Responsible Claims}
Evidence consistency does not establish truth, optimality, or hallucination
immunity. The verifier can inherit a misspecified outcome criterion, and
endogenous feedback complicates external validity. Small paired pilots support
mechanism diagnosis rather than population level significance or
backbone-independent claims. Prompt replay alone establishes action sensitivity;
utility language requires exact checkpoint continuation. Controlled fixed-rule
injection supports only a controlled-error response claim, not mitigation of
naturally generated proposal hallucinations.

\section{Conclusion}
EGRM reframes agent memory from storage and retrieval to an auditable question
of rule authority under feedback. The planned evaluation is designed to test
whether, and under which registered conditions, that lifecycle reduces exposure
to self-reinforcing errors.

\bibliographystyle{ACM-Reference-Format}
\bibliography{references}

\end{document}
